\section{Proof Overview}
\label{sec:overview}

In this section we give an overview of the proof of the inclusion $\RE \subseteq \MIP^*$. Since all interactive proof systems considered in the paper involve a single-round interaction between a classical verifier and two quantum provers sharing entanglement we generally use the language of nonlocal games to describe such proof systems, and often refer to the provers as ``players''. In a nonlocal game $\game$ (or simply ``game'' for short), the verifier can be described as the combination of two procedures: a \emph{question sampling} procedure that samples a pair of questions $(x,y)$ for the players according to a distribution $\mu$ (known to the verifier and the players), and a \emph{decision} procedure (also known to all parties) that takes as input the players' questions and their respective answers $a,b$ and evaluates a predicate $D(x,y,a,b)\in \{0,1\}$ to determine the verifier's acceptance or rejection. Given a description of a nonlocal game $\game$, recall that $\val^*(\game)$ denotes the \emph{entangled value} of the game, which is defined as the supremum~\eqref{eq:ent-val} of the players' success probability in the game over all finite-dimensional tensor product strategies. (We refer to Section~\ref{sec:games} for definitions regarding nonlocal games.) 

Our results establish the existence of transformations on \emph{families} of nonlocal games $\{\game_n\}_{n\in \N}$ having certain properties. In order to keep track of efficiency (and ultimately, computability) properties it is important to have a way to specify such families in a {uniform} manner. Towards this we introduce the following formalism. A \emph{uniformly generated family of games} is specified through a pair of Turing machines $\verifier = (\sampler,\decider)$ that satisfy certain conditions, in which case the pair is called a \emph{normal form verifier}. The Turing machine $\sampler$ (called a \emph{sampler}) takes as input an index $n \in \N$ and returns the description of a procedure that can be used to sample questions $(x,y)$ in the game (this procedure itself obeys a certain format associated with ``conditionally linear'' distributions, defined below). The Turing machine $\decider$ (called a \emph{decider}) takes as input an index $n$, questions $(x,y)$, and answers $(a,b)$, and returns a single-bit decision. For the sake of this proof overview we assume that the sampling and decision procedures run in time polynomial in the index $n$; we refer to the running time of these procedures as the \emph{complexity} of the verifier. 
	Given a normal form verifier $\verifier= (\sampler,\decider)$ we associate to it an infinite family of nonlocal games $\{\verifier_n \}$ indexed by positive integers in the natural way.  
	


The main technical result of this paper is a \emph{gap-preserving compression} transformation on normal form verifiers. The following theorem presents an informal summary of the properties of this transformation. Recall that for a  game $\game$ and probability $0\leq p \leq 1$, $\Ent(\game,p)$ denotes the minimum local dimension of an entangled state shared by the players in order for them to succeed in $\game$ with probability at least $p$.

\begin{theorem}[Gap-preserving compression of normal form verifiers, informal]
\label{thm:compression-informal}
There exists a polynomial-time Turing machine $\Compress$ that, when given as input the description of a normal form verifier $\verifier = (\sampler,\decider)$, outputs the description of another normal form verifier $\verifier' = (\sampler',\decider')$ that satisfies the following properties: for all $n \in \N$, letting $N = 2^n$,
\begin{enumerate}
\item (\textbf{Completeness}) If $\val^*(\verifier_N) = 1$ then $\val^*(\verifier'_n) = 1$.
\item (\textbf{Soundness}) If $\val^*(\verifier_N) \leq \frac{1}{2}$ then $\val^*(\verifier'_n) \leq \frac{1}{2}$.
\item (\textbf{Entanglement lower bound}) $\Ent(\verifier_n',\frac{1}{2}) \geq \max \{ \Ent(\verifier_N,\frac{1}{2}), 2^{2^{\Omega(n)}} \}$. 
\end{enumerate}
\end{theorem}
\noindent The formal version of this theorem is stated in Section~\ref{sec:compression} as Theorem~\ref{thm:compression}. The terminology \emph{compression} is motivated by the fact, implicit in the informal statement of the theorem, that the time complexity of the verifier's sampling and decision procedures in the game $\verifier_n'$, which is polynomial in $n$, is exponentially smaller than the time complexity of the verifier in the game $\verifier_N$, which is polynomial in $N$ and thus exponential in $n$. 

Before giving an overview of the proof of Theorem~\ref{thm:compression-informal} we sketch how the existence of a Turing machine $\Compress$ with the properties stated in the theorem implies the inclusion $\RE\subseteq \MIP^*$. Recall that the complexity class $\RE$ consists of all languages $L$ such that there is a Turing machine $\cal{M}$ that accepts instances $x$ in $L$, and does not accept instances $x$ that are not in $L$ (but is not required to terminate on such instances). To show $\RE \subseteq \MIP^*$ we give an $\MIP^*$ protocol for the Halting Problem, which is a complete problem for $\RE$. The Halting Problem is the language that consists of all Turing machine descriptions $\cal{M}$ such that $\cal{M}$ halts when run on an empty input tape. (For the purposes of this overview, we blur the distinction between a Turing machine and its description as a string of bits.) We give a procedure that given a Turing machine $\cal{M}$ as input returns the description of a normal form verifier $\verifier^{\cal{M}} = (\sampler^\machine,\decider^\machine)$ with the following properties. First, if $\cal{M}$ does eventually halt on an empty input tape, then it holds that for all $n \in \N$, $\val^*(\verifier^{\cal{M}}_n) = 1$. Second, if $\cal{M}$ does not halt then for all $n \in \N$, $\val^*(\verifier^{\cal{M}}_n) \leq \frac{1}{2}$. 

We describe the procedure that achieves this. Informally, the procedure returns the specification of a verifier $\verifier^{\cal{M}} = (\sampler^\machine,\decider^\machine)$ such that $\decider^\machine$ proceeds as follows: on input $(n,x,y,a,b)$ it first executes the Turing machine $\cal{M}$ for $n$ steps. If $\cal{M}$ halts, then $\decider^\machine$ accepts. Otherwise, $\decider^\machine$ computes the description of the compressed verifier $\verifier' = (\sampler',\decider')$ that is the output of $\Compress$ on input $\verifier^\machine$, then executes the decision procedure $\decider'(n,x,y,a,b)$ and accepts if and only if $\decider'$ accepts.\footnote{The fact that the decider $\decider^\machine$ can invoke the $\Compress$ procedure on itself follows from a well-known result in computability theory known as \emph{Kleene's recursion theorem} (also called \emph{Roger's fixed point theorem})~\cite{Kleene1954,Rogers1987}.} 

To show that this procedure achieves the claimed transformation, consider two cases. First, observe that if $\cal{M}$ eventually halts in some number of time steps $T$, then by definition $\val^*(\verifier^\machine_n) = 1$ for all $n\geq T$. Using Theorem~\ref{thm:compression-informal} along with an inductive argument it then follows that $\val^*(\verifier^\machine_n) = 1$ for all $n \geq 1$. Second, if $\cal{M}$ never halts, then observe that for any $n\geq 1$ Theorem~\ref{thm:compression-informal} implies two separate lower bounds on the amount of entanglement required to win the game $\verifier^\machine_n$ with probability at least $\frac{1}{2}$: the dimension is (a) at least $2^{2^{\Omega(n)}}$, and (b) at least the dimension needed to win the game $\verifier^\machine_{2^n}$ with probability at least $\frac{1}{2}$. Applying an inductive argument it follows that an \emph{infinite} amount of entanglement is needed to win the game $\verifier_n$ with any probability greater than $\frac{1}{2}$. Thus, a sequence of finite-dimension strategies for $\verifier_n$ cannot lead to a limiting value larger than $\frac{1}{2}$, and $\val^*(\verifier^\machine_n) \leq \frac{1}{2}$. 

We continue with an overview of the ideas behind the proof of Theorem~\ref{thm:compression-informal}.

\newcommand{\CompressNW}{\mathsf{Compress}^{\mathsf{NW}}}
 
\paragraph{Compression by introspection.}
To start, it is useful to review the protocol introduced in~\cite{NW19} to show
the inclusion $\NEEXP \subseteq \MIP^*$.
Fix an $\NEEXP$-complete language $L$.
The $\MIP^*$ protocol for $\NEXP$ from~\cite{natarajan2018two}, when scaled up
to decide languages from $\NEEXP$, yields a family of nonlocal games $\{ \game_z
\}$ that are indexed by instances $z \in \{0,1\}^*$.
The family of games decides $L$ in the sense that for all $z$, the game
$\game_z$ has entangled value $1$ if $z \in L$, and has entangled value at most
$\frac{1}{2}$ if $z \notin L$.
Furthermore, if $n = |z|$ is the length of $z$, the verifier of the game
$\game_z$ has complexity $\poly(N)=\exp(|z|)$ (recall that we use this
terminology to refer to an upper bound on the running time of the verifier's
sampling and decision procedure).
Thus, this family of games does not by itself yield an $\MIP^*$ protocol for
$L$.
To overcome this the main contribution in~\cite{NW19} is the design of an
efficient compression procedure $\CompressNW$ that applies specifically to the
family of games $\{\game_z \}$.
When given as input the description of $\game_z$, $\CompressNW$ returns the
description of a game $\game_z'$ such that if $\val^*(\game_z) = 1$, then
$\val^*(\game_z') = 1$, and if $\val^*(\game_z) \leq \frac{1}{2}$, then
$\val^*(\game_z') \leq \frac{1}{2}$.
Furthermore, the complexity of the verifier for $\game_z'$ is $\poly(n)$.
Thus the family of games $\{ \game_z' \}$ decides $L$ and this shows that
$\NEEXP \subseteq \MIP^*$, which is the best lower bound known on $\MIP^*$
prior to our work.
  

Presented in this way, it is natural to suggest iterating the procedure $\CompressNW$ to achieve e.g.\ the inclusion $\NEEEXP \subseteq \MIP^*$. To explain the difficulty in doing so, we give a little more detail on the compression procedure. It consists of two main steps: starting from $\game_z$, perform (1) \emph{question reduction}, and (2) \emph{answer reduction}. The goal of (1) is to reduce the length of the questions generated by the verifier in $\game_z$ from $\poly(N)$ to $\poly(n)$. The goal of (2) is to achieve the same with respect to the length of answers expected from the players. Furthermore, the complexity of the verifier of the resulting game $\game_z'$ should be reduced from $\poly(N)$ to  $\poly(n)$.
	
	Part (1) is achieved through a technique referred to as ``introspection'' where, rather than sampling questions $(x,y)$ of length $\poly(N)$ as in the game $\game_z$, the verifier instead executes a carefully crafted nonlocal game with the players that (a) requires questions of length $\poly(n)$, (b) checks that the players share $\poly(N)$ EPR pairs, and (c) checks that the players measure the EPR pairs in such a way as to \emph{sample for themselves} a question pair $(x,y)$ such that one player gets $x$ and the other player gets $y$. In other words, the players are essentially forced to introspectively ask themselves the questions of $\game_z$. 
	
	After question reduction, the players still respond with $\poly(N)$-length answers, which the verifier has to check satisfies the decision predicate of the original game $\game_z$. The goal of Part (2) is to enable the decision procedure to implement the verification procedure while not requiring the entire full-length answers from the players. In the answer reduction scheme of~\cite{NW19} this is achieved by having the verifier run a \emph{probabilistically checkable proof} (PCP) with the players so that they succinctly \emph{prove} that first, they have introspected questions $(x,y)$ from the correct distribution, and second, that they are able to generate $\poly(N)$-length answers $(a,b)$ that would satisfy the decision predicate of the original game $\game_z$ when executed on $(x,y)$ and $(a,b)$. Since the questions and answers in the PCP are of length $\poly(n)$, this achieves the desired answer length reduction.
	
	Iterating this scheme presents a number of immediate difficulties that have to do with the fact that the sampling and decision procedures of the verifier in $\game'_z$ do not have such a nice form as those in $\game_z$. First of all, the compression procedure of~\cite{NW19} can only ``introspect'' a specific question distribution of a nonlocal game from~\cite{natarajan2018two}; we call this distribution a ``low-degree test distribution''.\footnote{The nonlocal game of~\cite{natarajan2018two} is part of a more general family of games called ``low-degree tests'', which have been studied extensively for their applications in complexity theory~\cite{babai1991non}, property testing~\cite{rubinfeld1996robust}, and PCPs~\cite{arora1998probabilistic}. Generally, the question distribution of a low-degree test is a random question pair $(x,y)$ where $y$ is a randomly chosen affine subspace in $\F^m$ and $x$ is a uniformly random point on $y$, where $\F$ is a finite field and $m \geq 2$ is an integer. The specific low-degree test game in~\cite{natarajan2018two} (which is based on the low-degree test of~\cite{raz1997sub}) uses two-dimensional subspaces; thus the question distribution of~\cite{raz1997sub,natarajan2018two} is referred to as the ``plane-point distribution''.} However, the resulting question distribution of the question-reduced verifier, which is used to check the introspection, has a much more complex structure. A similar issue arises with the modifications required to perform answer reduction. In the PCP employed to achieve this the question distribution appears to be much more complex than the low-degree test distribution (this is in large part due to the need for a specially tailored PCP procedure that encodes separately different chunks of the witness, corresponding to answers from different players). As a result it is entirely unclear at first whether the question distribution used by the verifier in $\game'_z$ can be ``introspected'' for a second time. 
	
	To overcome these difficulties we identify a natural class of question distributions, called \emph{conditionally linear distributions}, that generalize the low-degree test distribution. We show that conditionally linear distributions can be ``introspected'' using conditionally linear distributions only, enabling recursive introspection. (In particular, they are a rich enough class to capture the types of question distributions produced by the compression scheme of~\cite{NW19}.) We define normal form verifiers by restricting their sampling procedure to generate conditionally linear question distributions, and this allows us to obtain the compression procedure on normal form verifiers described in Theorem~\ref{thm:compression-informal}. 
	
	Conceptually,  the identification of a natural class of distributions that is ``closed under introspection'' is a key step that enables the introspection technique to be applied recursively. (As we will see later, other closure properties of conditionally linear distributions, such as taking direct products, play an important role as well.) Since conditionally linear distributions are central to our construction we describe them next. 
	


\paragraph{Conditionally linear distributions.} 
Fix a vector space $V$ that is identified with $\F^m$, for a finite field $\F$ and integer $m$. Informally (see Definition~\ref{def:cl-func} for a precise definition), a function $L$ on $V$ is \emph{conditionally linear} (CL for short) if it can be evaluated by a procedure that takes the following form: (i) read a substring $z^{(1)}$ of $z$; (ii) evaluate a linear function $L_1$ on $z^{(1)}$; (iii) repeat steps (i) and (ii) with the remaining coordinates $z\backslash z^{(1)}$, such that the next steps are allowed to depend in an arbitrary way on $L_1(z^{(1)})$ but not directly on $z^{(1)}$ itself. What distinguishes a function of this form from an arbitrary function is that we restrict the number of iterations of (i)---(ii) to a constant number (at most $9$, in our case). (One may also think of CL functions as ``adaptively linear'' functions, where the number of ``levels'' of adaptivity is the number of iterations of (i)---(ii).) A distribution $\mu$ over pairs $(x,y) \in V\times V$ is called {conditionally linear}  if it is the image under a pair of {conditionally linear functions} $L^\alice,L^\bob: V\to V$ of the uniform distribution on $V$, i.e.\ $(x,y)\sim (L^\alice(z),L^\bob(z))$ for uniformly random $z\in V$. 

An important class of CL distributions are low-degree test distributions, which are distributions over question pairs $(x,y)$ where $y$ is a randomly chosen affine subspace of $\F^m$ and $x$ is a uniformly random point on $y$. We explain this for the case where the random subspace $y$ is one-dimensional (i.e.\ a line). Let $V = V_\xpt \oplus V_{\dir{}}$ where $V_\xpt = V_{\dir{}} = \F^m$. Let $L^\alice$ be the projection onto $V_\xpt$ (i.e. it maps $(x,v) \to (x,0)$ where $x \in V_\xpt$ and $v \in V_{\dir{}}$). Define $L^\bob: V \to V$ as the map $(x,v) \mapsto (L^\lnf_v(x),v)$ where $L^\lnf_v: V_\xpt \to V_\xpt$ is a linear map that, for every $v \in V_{\dir{}}$, projects onto a complementary subspace to the one-dimensional subspace of $V_\xpt$ spanned by $v$ (one can think of this as an ``orthogonal subspace'' to the span of $\{ v \}$). $L^\bob$ is conditionally linear because it can be seen as first reading the substring $v \in V_{\dir{}}$ (which can be interpreted as specifying the \emph{direction} of a line), and then applying a linear map $L^\lnf_v$ to $x \in V_\xpt$ (which can be interpreted as specifying a canonical point on the line $\ell = \{ x + tv : t \in \F \}$). It is not hard to see (and shown formally in Section~\ref{sec:ld-game}) that the distribution of $(L^\alice(z),L^\bob(z))$ for $z$ uniform in $V$, is identical (up to relabeling) to the low-degree test distribution $(x,\ell)$ where $\ell$ is a uniformly random affine line in $\F^m$, and $x$ is a uniformly random point on $\ell$.





Our main result about CL distributions, presented in Section~\ref{sec:introspection}, is that any CL distribution $\mu$, associated with a pair of CL functions $(L^\alice,L^\bob)$ over a linear space $V=F^m$, can be ``introspected'' using a CL distribution that is ``exponentially smaller'' than the initial distribution. Slightly more formally, to any CL distribution $\mu$ we associate a two-player game $\game_\mu$ (called the ``introspection game'') in which questions from the verifier are sampled from a CL distribution $\mu'$ over $\F^{m'}$ for some $m' = \poly\log(m)$ and such that in any successful strategy for the game $\game_\mu$, when the players are queried on a special question labeled $\Introspect$, they must respond with a pair $(x,y)$ that is approximately distributed according to $\mu$.
(The game allows us to do more: it allows us to conclude \emph{how} the players obtained $(x,y)$ --- by measuring shared EPR pairs in a specific basis --- and this will be important when using the game as part of a larger protocol that involves other checks.) 
Crucially for us, the distribution $\mu'$ only depends on a size parameter associated with $(L^\alice,L^\bob)$ (essentially, the integer $m$ together with the number of ``levels'' of adaptivity of $L^\alice$ and $L^\bob$), but not on any other structural property of $(L^\alice,L^\bob)$. Only the decision predicate for the introspection game $\game_\mu$ depends on the entire description of $(L^\alice,L^\bob)$.

We say a few words about the design of $\mu'$ and the associated introspection game, which borrow heavily from~\cite{NW19}. Building on the ``quantum low-degree test'' introduced in~\cite{natarajan2018low} it is already known how a verifier can force a pair of players to measure $m$ EPR pairs in either the computational or Hadamard basis and report the (necessarily identical) outcome $z$ obtained, all the while using questions of length polylogarithmic in $m$ only. The added difficulty is to ensure that a player obtains, and returns, precisely the information about $z$ that is contained in $L^\alice(z)$ (resp. $L^\bob(z)$), and not more. A simple example is the line-point distribution described earlier: there, the idea to ensure that e.g.\ the ``point'' player only obtains the first component, $x$ of $(x,v) \in V_\xpt \oplus V_{\dir{}}$, the verifier demands that the ``point'' player measures their qubits associated with the space $V_{\dir{}}$ in the Hadamard, instead of computational, basis; due to the uncertainty principle this has the effect of ``erasing'' the outcome in the computational basis.  The case of the ``line'' player is a little more complex: the goal is to ensure that, conditioned on the specification of the line $\ell$ received by the ``line'' player, the point $x$ received by the ``point'' player is uniformly random within $\ell$. This was shown to be possible in~\cite{NW19}.

We can now describe how samplers of normal form verifiers are defined: these are
Turing machines $\sampler$ that specify an infinite family of CL distributions
$\{ \mu_n \}$ by computing, for each index $n$, 
the CL functions $(L^{\alice,\, n},L^{\bob,\, n})$ associated with $\mu_n$. 
(See Definition~\ref{def:sampler} for a formal definition of samplers.)
Thus, the question distributions of a normal form verifier $\verifier =
(\sampler,\decider)$ are the CL distributions corresponding to $\sampler$.

    
    
\paragraph{Question reduction.} Just like the compression procedure of~\cite{NW19}, the compression procedure $\Compress$ of Theorem~\ref{thm:compression-informal} begins with performing question reduction on the input game. Given a normal form verifier $\verifier = (\sampler,\decider)$, the procedure $\Compress$ first computes a normal form verifier $\verifier^\intro = (\sampler^\intro,\decider^\intro)$ where for all $n \in \N$, the game $\verifier^\intro_n$ consists of playing the original game $\verifier_N$ where $N = 2^n$, except that instead of sampling the questions according to the CL distribution $\mu_N$ specified by the sampler $\sampler_N$, the verifier executes the introspection game $\game_{\mu_N}$  described in the previous subsection. Thus, in the game $\verifier^\intro_n$, when both players receive the question labeled $\Intro$ they are expected to sample $(x,y)$ respectively according to $\mu_N$, and respond with the sampled question together with answers $a,b$ respectively. The decider $\decider^\intro$ on index $n$ evaluates $\decider(N,x,y,a,b)$ and accepts if and only if $\decider$ accepts. As a result the time complexity of decider $\decider^\intro$ on index $n$ remains that of $\decider$, i.e.\ $\poly(N)$. However, the length of questions asked in $\verifier^\intro_n$ and the complexity of the sampler $\sampler^\intro$ are exponentially reduced, to $\poly(n)$.

For convenience we refer to the questions asked by the verifier in the ``question-reduced'' game $\verifier^\intro_n$ as ``small questions,'' and the questions that are introspected by the players in $\verifier^\intro_n$ (equivalently, the questions asked in the original game $\verifier_N$) as ``big questions.'' 

		
\paragraph{Answer reduction.} 
Having reduced the complexity of the question sampling, the next step in the compression procedure $\Compress$ is to reduce the complexity of decider $\decider^\intro$ from $\poly(N)$ to $\poly(n)$ (which necessarily implies reducing the answer length to $\poly(n)$). To achieve this the compression procedure computes a normal form verifier $\verifier^\ar = (\sampler^\ar,\decider^\ar)$ from $\verifier^\intro$ such that both the sampler and decider complexity in $\verifier^\ar$ are $\poly(n)$ (here, $\ar$ stands for ``answer reduction'').

Similarly to the answer reduction performed in~\cite{NW19}, at a high level this is achieved by composing the game $\verifier^\intro_n$ with a probabilistically checkable proof (PCP). In our context a PCP is a proof encoding that allows a verifier to check whether, given Turing machine $\cal{A}$ and time bound $T$ provided as input, there exists some input $x$ that $\cal{A}$ accepts in time $T$. The PCP proof can be computed from $\cal{A}$, $T$, and the accepting input (if it exists) and has length polynomial in $T$ and the description length $|\cal{A}|$ of $\cal{A}$. Crucially, the verifier can check a purported proof while only reading a constant number of symbols of it, each of length $\polylog(T,|\cal{A}|)$, and executing a verification procedure that runs in time $\polylog(T,|\cal{A}|)$. 

We use PCPs for answer reduction as follows. The verifier in the game $\verifier^\ar_n$ samples questions as $\verifier^\intro_n$ would and sends them to the players. Instead of receiving the introspected questions and answers $(x,y,a,b)$ for the original game $\verifier_N$ and running the decision procedure $\decider(N,x,y,a,b)$, the verifier instead asks the players to compute a PCP $\Pi$ for the statement that the original decider $\decider$ accepts the input $(N,x,y,a,b)$ in time $T=\poly(N)$. The verifier then samples additional questions for the players that ask them to return specific entries of the proof $\Pi$. Finally, upon receipt of the players' answers, the verifier executes the PCP verification procedure. Because of the efficiency of the PCP, both the sampling of the additional questions and the decision procedure can be executed in time $\poly(n)$.\footnote{This idea is inspired by the technique of composition in the PCP literature, in which the complexity of a verification procedure can be reduced by composing a proof system (often a PCP itself) with another PCP.}

This very rough sketch presents some immediate difficulties. A first difficulty is that in general no player by themselves has access to the entire input $(N,x,y,a,b)$ to $\decider$, so no player can compute the entire proof $\Pi$. We discuss this issue in the next paragraph. A second difficulty is that a black-box application of an existing PCP, as done in~\cite{NW19}, results in a question distribution for $\verifier_n^\ar$ (i.e.\ the sampling of the proof locations to be queried) that is rather complex --- and in particular, it may no longer fall within the framework of CL distributions for which we can do introspection. 
To avoid this, we design a bespoke PCP based on the classical MIP for NEXP (in particular, we borrow and adapt techniques from~\cite{ben2005simple,ben2006robust}). Two essential properties for us are that (i) the PCP proof is a collection of several low-degree polynomials,
    two of which are low-degree encodings of each player's big answer in the game $\verifier^\intro_n$, and (ii) verifying the proof only requires (a) running low-degree tests, (b) querying all polynomials at a uniformly random point, and (c) performing simple consistency checks. Property (i) allows us to eliminate the extra layer of encoding in~\cite{NW19}, who had to consider a PCP of proximity for a circuit applied to the
low-degree encodings of the players' big answers. Property (ii) allows us to ensure that the question distribution employed by $\verifier_n^\ar$ remains conditionally linear. 

\paragraph{Oracularization.}
The preceding paragraph raises a non-trivial difficulty. In order for the players to compute a proof for the claim that $\decider(N,x,y,a,b)=1$ they need to know the entire input $(x,y,a,b)$. However, in general a player only has access to their own question and answer: one player only knows $(x,a)$ and the other player knows $(y,b)$. The standard way of circumventing this difficulty is to consider an ``oracularized'' version of the game, where one player gets both questions $(x,y)$ and is able to determine both answers $(a,b)$, while the other player only gets one of the questions at random, and is only asked for one of the answers, that is then checked for consistency with the first player's answer. 

While this technique works well for games with classical players, when the players are allowed to use quantum strategies using entanglement oracularization does not, in general, preserve the completeness property of the game. To ensure that completeness is preserved we need an additional property of a completeness-achieving strategy for the original game: that there exists a \emph{commuting and consistent strategy} on all pairs of questions $(x,y)$ that are asked in the game with positive probability. Here commuting means that the measurement $\{A^x_a\}_a$ performed by the player receiving $x$ commutes with the measurement $\{B^y_b\}_b$ performed by the player receiving $y$.\footnote{We stress that the commuting property only applies to question pairs that occur with positive probability, and does not mean that \emph{all} pairs of measurement operators are required to commute; indeed this would imply that the strategy is effectively classical.} Consistent means that if both players perform measurements associated with the same question they obtain the same answer. If both properties hold then in the oracularized game when one player receives a pair $(x,y)$ and the other player receives the question $x$ (say), the first player can simultaneously measure both $\{A^x_a\}_a$ and $\{B^y_b\}_b$ on their own space to obtain a pair of answers $(a,b)$, and the second player can measure $\{A^x_a\}_a$ to obtain a consistent answer $a$.

For answer reduction to be possible it is thus applied to the {oracularized} version of the introspection game $\verifier^\intro_n$. This in turn requires us to ensure that the introspection game $\verifier^\intro_n$ has a commuting and consistent strategy achieving value $1$ whenever it is the case that $\val^*(\verifier_n^\intro)=1$. For this property to hold we verify that it holds for the initial game that is used to seed the compression procedure (this is true because we can start with an $\MIP^*$ protocol for $\NEXP$ for which there exists a perfect classical strategy) and we also ensure that each of the transformations of the compression protocol (question reduction, answer reduction, and parallel repetition described next) maintains it.


\paragraph{Parallel repetition.}
The combined steps of question reduction (via introspection) and answer reduction (via PCP composition) result in a game $\verifier_n^\ar$ such that the complexity of the verifier is $\poly(n)$. Furthermore, if the original game $\verifier_N$ has value $1$, then $\verifier_n^\ar$ also has value $1$. Unfortunately the sequence of transformations incurs a loss in the soundness parameters: if $\val^*(\verifier_N) \leq \frac{1}{2}$, then we can only establish that $\val^*(\verifier^\ar_n) \leq 1 - C$ for some positive constant $C < \frac{1}{2}$ (we call $C$ the \emph{soundness gap}). Such a loss would prevent us from recursively applying the compression procedure $\Compress$ an arbitrary number of times, which is needed to obtain the desired complexity results for $\MIP^*$. 

To overcome this we need a final transformation to restore the soundness gap of the games after answer reduction to a constant larger than $\frac{1}{2}$. To achieve this we use the technique of {parallel repetition}. The parallel repetition of a game $\game$ is another nonlocal game $\game^k$, for some number of repetitions $k$, which consists of playing $k$ independent and simultaneous instances of $\game$ and accepting if and only if all $k$ instances accept. Intuitively, parallel repetition is meant to decrease the value of a game $\game$ exponentially fast in $k$, provided $\val^*(\game) < 1$ to begin with. However, it is an open question of whether this is generally true for the entangled value $\val^*$. 

Nevertheless, some variants of parallel repetition are known to achieve exponential amplification. We use  a variant  called ``anchored parallel repetition'' and introduced in~\cite{bavarian2017hardness}. This allows us to devise a transformation that efficiently amplifies the soundness gap to a constant. The resulting game $\verifier^\rep_n$ has the property that if $\val^*(\verifier^\ar_n) = 1$, then $\val^*(\verifier^\rep_n) = 1$ (and moreover this is achieved using a commuting and consistent strategy), whereas if $\val^*(\verifier^\ar_n) \leq 1 - C$ for some universal constant $C > 0$ then $\val^*(\verifier^\rep_n) \leq \frac{1}{2}$. Furthermore, we have the additional property, essential for us, that good strategies in the game $\verifier^\rep_n$ require as much entanglement as good strategies in the game $\verifier^\ar_n$ (which in turn require as much entanglement as good strategies in the game $\verifier_N$). The complexity of the verifier in $\verifier^\rep_n$ remains $\poly(n)$.

The anchored parallel repetition procedure, when applied to a normal form verifier, also yields a normal form verifier: this is because the direct product of CL distributions is still conditionally linear. 


\paragraph{Putting it all together.} This completes the overview of the transformations performed by the compression procedure $\Compress$ of Theorem~\ref{thm:compression-informal}. To summarize, given an input normal form verifier $\verifier$, question reduction is applied to obtain $\verifier^\intro$, answer reduction is applied to the oracularized version of $\verifier^\intro$ to obtain $\verifier^\ar$, and  anchored parallel repetition is applied to obtain $\verifier^\rep$, which is returned by the compression procedure. Each of these transformations preserves completeness (including the commuting and consistent properties of a value-1 strategy) as well as the entanglement requirements of each game; moreover, the overall transformation preserves soundness. 
 
\section{Preliminaries}
